\documentclass[11pt,a4paper]{article}
\usepackage[T1]{fontenc}
\usepackage[utf8]{inputenc}
\usepackage{lmodern}
\usepackage{amsmath,amssymb,amsthm}
\usepackage[margin=1.15in]{geometry}
\usepackage[hidelinks]{hyperref}
\emergencystretch=3em

\newtheorem{theorem}{Theorem}
\renewcommand{\thetheorem}{\Alph{theorem}}
\newtheorem{lemma}{Lemma}[section]
\newtheorem{corollary}[theorem]{Corollary}
\theoremstyle{definition}
\newtheorem{definition}[lemma]{Definition}
\newtheorem{remark}[lemma]{Remark}

\newcommand{\Z}{\mathbb{Z}}
\newcommand{\Q}{\mathbb{Q}}
\newcommand{\C}{\mathbb{C}}
\renewcommand{\IJ}{I_J}
\newcommand{\sJ}{s_J}
\newcommand{\Pic}{\operatorname{Pic}}
\newcommand{\Hom}{\operatorname{Hom}}
\newcommand{\rad}{\operatorname{rad}}
\newcommand{\cl}{\operatorname{cl}}

\title{Multiplication by $2$ and by $6$ through products of Jacobians,
and the Jacobian index of a very general abelian variety of dimension
$4$, $5$ and $6$\thanks{AI-produced. Theorem~\ref{thm:A} is due to a run of
GPT-6 Astra (OpenAI, ChatGPT Pro mode); Theorem~\ref{thm:B} and parts (i)--(ii)
of Corollary~\ref{cor:C} to Claude Opus~5.5 (Anthropic) at effort max; the
text, Theorem~\ref{thm:D} and part (iii) of Corollary~\ref{cor:C} to Claude
Opus~5.5 at effort max; see Section~\ref{sec:provenance}. No human has
checked this text. To the extent that any rights exist in it, they are
dedicated to the public domain under CC0~1.0.}}
\author{}
\date{25 November 2026}

\begin{document}
\maketitle

\begin{abstract}
For an abelian variety $A$ over an algebraically closed field, let
$\sJ(A)$ be the least positive integer $n$ such that multiplication by $n$
factors through a product of Jacobians of curves; $\sJ(A)=1$ exactly when
$A$ is a direct summand of such a product. We prove, in every
characteristic, that multiplication by $2$ factors through a product of
Jacobians on every principally polarized abelian variety of dimension at
most $5$, and multiplication by $6$ on every one of dimension $6$; that
$\sJ(A)$ divides $\gcd(2,\delta(A))$ for every abelian surface and
$\gcd(6,\delta(A))$ for every abelian threefold, where $\delta(A)$ is the
greatest common divisor of the Euler characteristics of line bundles on
$A$. The known bound is $(2g-1)!$. For a complex principally polarized $A$
with minimal class $\gamma=\theta^{g-1}/(g-1)!$, let $m(A)$ be the least
positive multiple of $\gamma$ that is the class of an algebraic cycle;
Matsusaka's construction gives $\sJ(A)\mid m(A)$, and we prove
$m(A)\mid\sJ(A)^2$. With the theorems of
Engel, de Gaay Fortman and Schreieder this gives the exact values
$\sJ(A)=2$ for a very general principally polarized abelian variety of
dimension $4$ and $5$, and $\sJ(A)=6$ in dimension $6$; and every prime
$p\le g/2$ divides $\sJ(A)$ in dimension $g$. Finally, $\sJ(A)$ is exactly
the least integer $n$ for which $n$ times the inverse of the isomorphism
$H^{2g-1}(A)(g-1)\to H^1(\widehat A)$ is induced by an integral algebraic
cycle; this answers, for abelian varieties, the second question that
Grothendieck asks on sheet 70 of cote~14 of his Montpellier archive. The
results come from a run of GPT-6 Astra (ChatGPT Pro mode) and from Claude
Opus~5.5 at effort max; they were checked and the text written by Claude
Opus~5.5 at effort max.
\end{abstract}

\section{Results}\label{sec:results}

Throughout, $k$ is an algebraically closed field, curves are smooth,
projective and connected, and a \emph{product of Jacobians} is
$J=J(C_1)\times\dots\times J(C_r)$.

\begin{definition}\label{def}
For an abelian variety $A$ over $k$, let $\IJ(A)$ be the set of $n\in\Z$
such that $[n]_A=\rho\circ\iota$ for some product of Jacobians $J$ and
homomorphisms $\iota\colon A\to J$, $\rho\colon J\to A$. It is an ideal
(use $J\times J'$ for sums) and is non-zero, since $A$ is an isogeny factor
of a Jacobian; $\sJ(A)$ denotes its positive generator. Let
$\delta(A)=\gcd\{\chi(L):L\in\Pic(A)\}$. For a principally polarized
complex $(A,\theta)$ of dimension $g$, put $\gamma=\theta^{g-1}/(g-1)!\in
H^{2g-2}(A,\Z)$, let $M(A)$ be the set of $n$ such that $n\gamma$ is the
class of an integral algebraic $1$-cycle, and let $m(A)$ be its positive
generator; $(g-1)!\in M(A)$.
\end{definition}

\begin{theorem}\label{thm:A}
Let $k$ have any characteristic.
\begin{enumerate}
\item[(a)] If $B$ is a principally polarized abelian variety of dimension at
most $5$, then $2\in\IJ(B)$; if it has dimension $6$, then $6\in\IJ(B)$.
\item[(b)] If $A$ is an abelian surface, then $\sJ(A)\mid\gcd(2,\delta(A))$.
If $A$ is an abelian threefold, then $\sJ(A)\mid\gcd(6,\delta(A))$.
\end{enumerate}
In particular an abelian surface with $\delta(A)$ odd, and an abelian
threefold with $\gcd(\delta(A),6)=1$, is a direct summand of a product of
Jacobians.
\end{theorem}

\begin{theorem}\label{thm:B}
For every complex principally polarized abelian variety $A$,
\[
\sJ(A)\mid m(A)\mid \sJ(A)^2 .
\]
Hence $\rad\sJ(A)=\rad m(A)$, and $\sJ(A)=m(A)$ whenever $m(A)$ is
squarefree.
\end{theorem}

\begin{corollary}\label{cor:C}
Let $A$ be a very general complex principally polarized abelian variety of
dimension $g$.
\begin{enumerate}
\item[(i)] If $g\in\{4,5\}$, then $\sJ(A)=2$.
\item[(ii)] If $g=6$, then $\sJ(A)=6$.
\item[(iii)] If $g\ge4$, every prime $p\le g/2$ divides $\sJ(A)$, and
$\sJ(A)\mid m(A)\mid (g-1)!$.
\end{enumerate}
\end{corollary}

For an abelian variety $A$ of dimension $g$ and a prime $\ell\ne\operatorname{char}k$,
let
\[
\psi_A\colon H^{2g-1}(A,\Z_\ell)(g-1)\xrightarrow{\ \sim\ }H^1(\widehat A,\Z_\ell)
\]
be the canonical isomorphism: Poincar\'e duality
$H^{2g-1}(A,\Z_\ell)(g)\cong T_\ell A$, followed by
$T_\ell A(-1)\cong H^1(\widehat A,\Z_\ell)$ given by the Weil pairing
(over $\C$ one may use Betti cohomology with $\Z$-coefficients). It is
defined up to a universal sign, which plays no role below. An integral
algebraic $1$-cycle $Z$ on $\widehat A\times A$ acts as a correspondence
$Z_*\colon H^1(\widehat A,\Z_\ell)\to H^{2g-1}(A,\Z_\ell)(g-1)$.

\begin{theorem}\label{thm:D}
For $n\in\Z$ the following are equivalent: (i) $n\,\psi_A^{-1}=Z_*$ for an
integral algebraic $1$-cycle $Z$ on $\widehat A\times A$; (ii)
$n\in\IJ(A)$. In particular $\sJ(A)$ is the least positive integer $n$ such
that $n\,\psi_A^{-1}$ is induced by an integral algebraic cycle.
\end{theorem}

\paragraph{Why it matters.} On sheet~70 of cote~14 of Grothendieck's
Montpellier archive (scan~71), after constructing rational inverses of the
Lefschetz-type maps by $1$-cycles (Prop.~2.7, scans~68--70), Grothendieck
asks two questions. The second reads: is there always an \emph{integral}
cycle giving the inverse of $H^{2n-1}(X)(n-1)\to H^1(B_X)$, $B_X$ the
Picard variety of $X$? He notes that if the first question has a positive
answer, it suffices to answer the second for abelian varieties.
Theorem~\ref{thm:D} identifies the obstruction for an abelian variety with
$\sJ(A)$; Theorem~\ref{thm:A} and Corollary~\ref{cor:C} compute it. The
answer is positive for all abelian surfaces with $\delta$ odd, all
threefolds with $\delta$ prime to $6$ and all principally polarized ones of
dimension at most $3$, and negative for a very general principally
polarized abelian variety of dimension at least $4$, where it fails by
exactly $2$ in dimensions $4$, $5$ and by exactly $6$ in dimension~$6$.
By the criterion of Voisin \cite[Thm.~1.4]{Vo}, as recalled in
\cite[\S1.2]{EGFS}, $\sJ=1$ for all abelian varieties is equivalent to the
integral Hodge conjecture for curve classes on all abelian varieties, which
Engel, de Gaay Fortman and Schreieder disproved \cite{EGFS}; $\sJ(A)$
measures how far a given $A$ is from that splitting.

\paragraph{Credit.} Theorem~\ref{thm:A} is due to a run of GPT-6 Astra
(OpenAI, ChatGPT Pro mode) of this project; Theorem~\ref{thm:B} and
Corollary~\ref{cor:C}\,(i)--(ii) to Claude Opus~5.5 (Anthropic) at effort
max; Theorem~\ref{thm:D} and Corollary~\ref{cor:C}\,(iii) to the writer,
Claude Opus~5.5 at effort max. Details are in Section~\ref{sec:provenance}.

\paragraph{Context: what is known and what is new.} Known: the implication
\emph{$n\gamma$ algebraic $\Rightarrow$ $n\in\IJ(A)$} for principally
polarized $A$ over any field, by Matsusaka's construction
\cite[Prop.~3.4, eqs.~(3.1)--(3.3)]{SS}, which applied to the unimodular
Poincar\'e class of $A\times\widehat A$ gives $\sJ(A)\mid(2g-1)!$ (see also
\cite[\S3]{BdGF}); the
case $\sJ=1\Leftrightarrow m=1$ \cite[Thm.~1.1]{SS}, over $\C$ due to
Beckmann and de Gaay Fortman \cite{BdGF} and Voisin \cite{Vo}; the integral
Hodge conjecture for $1$-cycles on products of Jacobians \cite[Thm.~1.2]{BdGF};
the values $m(A)=2$ for very general $A$ of dimension $4,5$
\cite[Cor.~1.2]{EGFS}, $m(A)=6$ in dimension $6$ \cite[Thm.~1.2]{EGFS2},
the divisibility of $m(A)$ by every prime $p\le g/2$ \cite[Thm.~1.1,
Cor.~1.2]{ES}, and the fact that a very general principally polarized $A$
of dimension $\ge4$ is not a direct summand of a product of Jacobians and
admits no isogeny of odd degree $A\times Y\to J$ \cite[Cor.~1.5]{EGFS}.
These already give Corollary~\ref{cor:C}\,(i): $\sJ(A)\mid m(A)=2$ by the
implication above and \cite[Cor.~1.2]{EGFS}, and $\sJ(A)\ne1$ by
\cite[Cor.~1.5]{EGFS} (Remark~\ref{rem:even}). New, to our knowledge after the
search described in Section~\ref{sec:provenance}: Theorem~\ref{thm:A} (for
surfaces the known bound is $3!=6$, for threefolds $5!=120$), the
divisibility $m(A)\mid\sJ(A)^2$ of Theorem~\ref{thm:B}, which parts (ii) and
(iii) of Corollary~\ref{cor:C} need, and Theorem~\ref{thm:D}.

\paragraph{Scope.} Theorem~\ref{thm:B} and Corollary~\ref{cor:C} are over
$\C$; we make no claim on very general abelian varieties in positive
characteristic. Theorem~\ref{thm:A} leaves open the abelian surfaces with
$\delta(A)$ even, e.g.\ a surface of type $(1,2)$ with
$\operatorname{End}(A)=\Z$, where $\sJ(A)\in\{1,2\}$, and the threefolds
with $\gcd(\delta(A),6)>1$, e.g.\ of type $(1,1,3)$ or $(1,2,2)$. The first
question of sheet~70 is not treated.

\paragraph{Proof overview.} Theorem~\ref{thm:A}(a) uses Beauville's
generalized Prym varieties in dimension $\le5$ and the Prym--Tyurin
uniformization of Alexeev, Donagi, Farkas, Izadi and Ortega in dimension
$6$, which give curve classes $2\gamma$ and $6\gamma$; lifting to
characteristic zero and specializing these cycles gives every
characteristic, and Matsusaka's construction turns them into factorizations.
Part (b) applies (a) to the principally polarized variety
$A\times\widehat A$ after Mumford's descent of a polarization to a principal
one. Theorem~\ref{thm:B} pushes the minimal class into a product of
Jacobians, where it is algebraic by \cite{BdGF}, and back. Theorem~\ref{thm:D}
reads a $1$-cycle on $\widehat A\times A$ as a sum of homomorphisms through
Jacobians.

\section{Proofs}\label{sec:proofs}

\begin{lemma}[Matsusaka; {\cite[Prop.~3.4]{SS}}]\label{lem:mats}
Let $(B,\theta)$ be principally polarized of dimension $h$ over $k$. If an
integral $1$-cycle on $B$ has numerical class $n\,\theta^{h-1}/(h-1)!$,
then $n\in\IJ(B)$.
\end{lemma}

This is \cite[Prop.~3.4]{SS} with $\deg(\theta^h/h!)=1$: the composite
$B\to\widehat B\to\prod J(\widetilde C_i)\to B$ built from the normalized
components $\widetilde C_i$ of the cycle is $[n]$ by
\cite[eqs.~(3.1)--(3.3)]{SS}.

\begin{lemma}\label{lem:delta}
(a) $\delta(A)=\gcd\{\chi(L):L\text{ ample}\}$. (b) If $L$ is ample and
$n\in\IJ(B)$ for every principally polarized $B$ of dimension $\dim A$, then
$n\,\chi(L)\in\IJ(A)$; hence $n\,\delta(A)\in\IJ(A)$. (c) If
$\dim A\le3$, then $\delta(A)\in\IJ(A)$.
\end{lemma}

\begin{proof}
(a) Let $g=\dim A$, $L\in\Pic(A)$ and $H$ ample. By Riemann--Roch
$p(t)=\chi(L\otimes H^{t})$ is a polynomial of degree $\le g$ with integer
values, so for every $T\ge1$ the $(g+1)$-st difference of $s\mapsto p(sT)$
vanishes at $0$:
$\chi(L)=\sum_{j=1}^{g+1}(-1)^{j+1}\binom{g+1}{j}\,\chi(L\otimes H^{jT})$.
For $T$ large all $L\otimes H^{jT}$, $j\ge1$, are ample.

(b) By Mumford's descent theorem in its group-scheme form
\cite[\S23, Thm.~4 and Cor.~1, pp.~233--234]{Mu}, there are a maximal
isotropic subgroup scheme $H\subset K(L)$, the quotient isogeny $f\colon
A\to B=A/H$ and a principal polarization $M$ on $B$ with $f^*M\cong L$;
hence $\deg f=\chi(L)$, and there is $f'\colon B\to A$ with
$f'f=[\chi(L)]$, since $\ker f\subset A[\deg f]$. If $[n]_B=\rho\iota$
through $J$, then $[n\chi(L)]_A=f'[n]_Bf=(f'\rho)(\iota f)$. Now use (a).

(c) Principally polarized abelian varieties of dimension $\le3$ are
products of Jacobians \cite[Thm.~4]{OU}, so (b) applies with $n=1$.
\end{proof}

\begin{proof}[Proof of Theorem~\ref{thm:A}]
(a) First let $\operatorname{char}k=0$; by the Lefschetz principle we may
take $k=\C$. If $\dim B=h\le5$, by Beauville \cite[Prop.~2.2,
Thm.~5.4, Thm.~6.5(i)]{Be} either $B$ is a product of Jacobians or there are
a product of Jacobians $J$, with principal polarization $\lambda_J$, and a
homomorphism $j\colon B\to J$ with $j^*\Theta_J\equiv2\Xi$, $\Xi$ the
polarization of $B$. Put $q=\lambda_B^{-1}j^\vee\lambda_J\colon J\to B$; then
$qj=\lambda_B^{-1}\varphi_{j^*\Theta_J}=[2]$, so $2\in\IJ(B)$. Pushing the
Abel--Jacobi curves of the factors of $J$ forward by $q$ gives an integral
$1$-cycle whose homomorphism $\widehat B\to B$ is
$q\lambda_J^{-1}q^\vee=\lambda_B^{-1}(j^\vee\lambda_Jj)\lambda_B^{-1}
=2\lambda_B^{-1}$, so its class is $2\,\Xi^{h-1}/(h-1)!$. If $\dim B=6$,
Alexeev, Donagi, Farkas, Izadi and Ortega \cite[Cor.~2, p.~165]{ADFIO}
give an effective curve of class $6\,\Xi^5/5!$, and Lemma~\ref{lem:mats}
gives $6\in\IJ(B)$.

Now let $\operatorname{char}k=p>0$. A principal polarization is separable,
so $(B,\Xi)$ lifts to a principally polarized abelian scheme over $W(k)$
\cite[Cor.~2.4.2 and its proof, p.~287]{Oo}. On its geometric generic fibre
the class $c\,\Xi^{h-1}/(h-1)!$, $c=2$ for $h\le5$ and $c=6$ for $h=6$, is
algebraic by the case just treated. The cycle and its rational equivalences
are defined over a finite extension; after the corresponding finite
extension of discrete valuation rings, its closure specializes to an
integral $1$-cycle on $B$ \cite[Ex.~20.3.5, p.~400]{Fu}, with the same class
in $\ell$-adic cohomology by smooth and proper base change. Lemma~\ref{lem:mats}
gives $c\in\IJ(B)$.

(b) Let $g=\dim A\in\{2,3\}$ and $X=A\times\widehat A$, of dimension $2g$,
with the Poincar\'e bundle $P$. Its homomorphism
$\varphi_P\colon X\to\widehat X=\widehat A\times A$ is $(a,\xi)\mapsto
(\xi,a)$ up to sign, an isomorphism, so $\chi(P)^2=\deg\varphi_P=1$ by the
Riemann--Roch theorem \cite[\S16]{Mu}, and $\delta(X)=1$. By (a), $c\in\IJ(B)$
for every principally polarized $B$ of dimension $2g$, with $c=2$ for $g=2$
and $c=6$ for $g=3$. Lemma~\ref{lem:delta}(b) gives
$c=c\,\delta(X)\in\IJ(X)$. Since $A$ is a direct factor of $X$, composing
with $A\to X\to A$ gives $c\in\IJ(A)$. With Lemma~\ref{lem:delta}(c),
$\gcd(c,\delta(A))\in\IJ(A)$.
\end{proof}

\begin{proof}[Proof of Theorem~\ref{thm:B}]
Lemma~\ref{lem:mats} gives $M(A)\subset\IJ(A)$, so $\sJ\mid m$. Let
$n=\sJ(A)$ and $[n]=\rho\circ\iota$ with $\iota\colon A\to J$. The class
$\iota_*\gamma\in H_2(J,\Z)$ is an integral Hodge class, because $\iota_*$ is
a morphism of Hodge structures. By Beckmann and de Gaay Fortman
\cite[Thm.~1.2]{BdGF} the integral Hodge conjecture for $1$-cycles holds on
products of Jacobians, so $\iota_*\gamma=[Z]$ for an integral $1$-cycle $Z$
on $J$. Then $\rho_*[Z]=[n]_*\gamma=n^2\gamma$, because $[n]$ acts on
$H_2(A,\Z)=\bigwedge^2H_1(A,\Z)$ by $n^2$. Hence $n^2\in M(A)$ and
$m\mid\sJ^2$. The two divisibilities give $\rad\sJ=\rad m$; if $m$ is
squarefree, $m\mid\sJ^2$ forces $m\mid\sJ$, so $\sJ=m$.
\end{proof}

\begin{proof}[Proof of Corollary~\ref{cor:C}]
For very general $A$, $m(A)$ is the integer $N_g$ of
\cite[Question~1.1]{EGFS2}. (i) By \cite[Cor.~1.2]{EGFS} with $c=g-1$, the
algebraic classes in $H^{2g-2}(A,\Z)$ are $2\Z\gamma$, so $m(A)=2$. (ii)
By \cite[Thm.~1.2]{EGFS2}, $m(A)=6$. Both are squarefree;
Theorem~\ref{thm:B} gives $\sJ=m$. (iii) By \cite[Thm.~1.1 and
Cor.~1.2]{ES} with $k=g-1$, every prime $p\le g/2$ divides $m(A)$, hence
$\sJ(A)$ by Theorem~\ref{thm:B}; and $\sJ\mid m\mid(g-1)!$.
\end{proof}

\begin{remark}\label{rem:even}
That $\sJ(A)$ is even for $g\ge4$ follows also from \cite[Cor.~1.5]{EGFS}
without Theorem~\ref{thm:B}: if $vu=[s]$ with $u\colon A\to J$, $v\colon
J\to A$ and $s$ odd, then $(x,y)\mapsto u(x)+y$ is an isogeny
$A\times(\ker v)^0\to J$ whose kernel embeds into $A[s]$, so its degree is
odd. With $\sJ(A)\mid m(A)=2$ \cite[Cor.~1.2]{EGFS} this gives
Corollary~\ref{cor:C}\,(i) without Theorem~\ref{thm:B}. For $g=6$ it leaves
$\sJ(A)\in\{2,6\}$; that $3\mid\sJ(A)$ needs Theorem~\ref{thm:B}.
\end{remark}

\section{Grothendieck's second question}\label{sec:D}

\begin{proof}[Proof of Theorem~\ref{thm:D}]
Let $C$ be a curve, $f=(p,q)\colon C\to\widehat A\times A$, and $Z=f_*[C]$.
The correspondence $Z_*$ is $q_*p^*$, which factors through $H^1(C,\Z_\ell)$.
Under $H^1(C,\Z_\ell)\cong T_\ell J(C)(-1)$, the Gysin map
$q_*\colon H^1(C)\to H^{2g-1}(A)(g-1)$ becomes $T_\ell$ of the Albanese map
$q_*\colon J(C)\to A$, and $p^*\colon H^1(\widehat A)\to H^1(C)$ becomes,
up to the autoduality sign of $J(C)$, $T_\ell$ of the pull-back
$p^*\colon A=\Pic^0(\widehat A)\to\Pic^0(C)=J(C)$. Hence, up to a universal
sign, $\psi_A\circ Z_*=T_\ell(\varphi_Z)$ with $\varphi_Z=q_*\circ p^*\colon
A\to J(C)\to A$, and this extends additively to integral $1$-cycles
$Z=\sum n_if_{i*}[\widetilde C_i]$, $\varphi_Z=\sum n_i\,q_{i*}p_i^*$.

(i)$\Rightarrow$(ii): if $Z_*=n\psi_A^{-1}$, then $T_\ell(\varphi_Z)=\pm n$,
so $\varphi_Z=\pm[n]$ because $\Hom(A,A)\to\operatorname{End}T_\ell A$ is
injective; $\varphi_Z$ factors through $\prod J(\widetilde C_i)$, so
$n\in\IJ(A)$.

(ii)$\Rightarrow$(i): let $[n]=\rho\iota$ with $\iota=(\iota_a)\colon A\to
\prod J(C_a)$ and $\rho=\sum\rho_a$. Let $\alpha_a\colon C_a\to J(C_a)$ be an
Abel--Jacobi map, $\lambda_a$ the principal polarization of $J(C_a)$ and
$\widehat\iota_a\colon\widehat{J(C_a)}\to\widehat A$ the dual of $\iota_a$.
For $f_a=(\widehat\iota_a\lambda_a\alpha_a,\rho_a\alpha_a)\colon C_a\to
\widehat A\times A$ one finds $p_a^*=\alpha_a^*\lambda_a\iota_a$ and
$q_{a*}=\rho_a\alpha_{a*}=\rho_a$, and $\alpha_a^*=-\lambda_a^{-1}$ on
$\Pic^0(J(C_a))$ \cite[Lemma~6.9]{Mi}. Hence
$\varphi_{Z}=-\sum_a\rho_a\iota_a=-[n]$ for $Z=\sum_af_{a*}[C_a]$, and
$\pm Z$ induces $n\psi_A^{-1}$.
\end{proof}

\section{Provenance and checks}\label{sec:provenance}

This paper is AI-produced. No human has proved, written or checked any part
of it.

\paragraph{What is whose.} Theorem~\ref{thm:A} is due to a run of GPT-6
Astra (OpenAI, ChatGPT Pro mode) of this project (run RESEARCH\_SA06, ended
19 November 2026; its statements C.4j, C.4m--C.4p), in answer to a research
prompt on the integral questions of cote~14 of Grothendieck's archive.
Theorem~\ref{thm:B} and Corollary~\ref{cor:C}\,(i)--(ii) are due to Claude
Opus~5.5 at effort max (agent a5d5ce46, 24 November 2026), in a map of the
project's geometry results. Theorem~\ref{thm:D}, Corollary~\ref{cor:C}\,(iii),
Remark~\ref{rem:even} and the text are due to the writer, Claude Opus~5.5 at
effort max (agent a0012f1f, 25 November 2026).

\paragraph{Checks.} The run's own audits (GPT-6 Astra, ChatGPT Pro
mode) passed. The writer re-proved Theorem~\ref{thm:A}(b) and the Matsusaka
coefficient independently on 22 November 2026, with exact computations, read
the Prym step of Theorem~\ref{thm:A}(a) in full and checked its sixfold step
at the intake of the run;
it read Theorem~\ref{thm:B} and Corollary~\ref{cor:C}\,(i)--(ii) a second
time on 25 November 2026, with the cited statements opened. On 25 November
2026 the main session of the project (Claude Opus 5.5 at effort max) read every
page and checked Theorem~\ref{thm:B}, Corollary~\ref{cor:C},
Remark~\ref{rem:even} and Theorem~\ref{thm:D} line by line.

\paragraph{The search of the literature.} On 25 November 2026 the writer
read the statements cited from \cite{BdGF,EGFS,EGFS2,ES,SS} and searched
arXiv (all papers of de Gaay Fortman, including \cite{GFS}, and the
abstracts mentioning products of Jacobians), the works citing
\cite{BdGF,EGFS,EGFS2,SS} listed by OpenAlex, and the Hexagon catalogue. No
statement of Theorems~\ref{thm:A}, \ref{thm:B}, \ref{thm:D} or of the exact
values in Corollary~\ref{cor:C} was found. The writer's search does not cover
full texts; the main session then had the full texts of \cite{SS} and
\cite{EGFS2}, and of \cite{EGFS} through its Section~3, searched for the
statements claimed here as new (by a fetch tool's reading model; none is
there), and opened \cite{ES} at its abstract. The locators of \cite{ADFIO,Be,Fu,Mi,Mu,OU,Oo,Vo} were opened by the
run, not by the writer. The full record is in the file PROVENANCE.md of the
paper's workbench.

\begin{thebibliography}{EGFS2}
\bibitem[ADFIO]{ADFIO} V. Alexeev, R. Donagi, G. Farkas, E. Izadi, A. Ortega,
\emph{The uniformization of the moduli space of principally polarized abelian
$6$-folds}, J. Reine Angew. Math. 761 (2020), 163--217.
\bibitem[Be]{Be} A. Beauville, \emph{Prym varieties and the Schottky
problem}, Invent. Math. 41 (1977), 149--196.
\bibitem[BdGF]{BdGF} T. Beckmann, O. de Gaay Fortman, \emph{Integral Fourier
transforms and the integral Hodge conjecture for one-cycles on abelian
varieties}, Compos. Math. 159 (2023), 1188--1213; arXiv:2202.05230v2.
\bibitem[EGFS]{EGFS} P. Engel, O. de Gaay Fortman, S. Schreieder,
\emph{Matroids and the integral Hodge conjecture for abelian varieties},
arXiv:2507.15704v3 (2026).
\bibitem[EGFS2]{EGFS2} P. Engel, O. de Gaay Fortman, S. Schreieder,
\emph{Optimality of the Prym--Tyurin construction for $\mathcal A_6$},
arXiv:2512.04902v1 (2025).
\bibitem[ES]{ES} P. Engel, S. Schreieder, \emph{On the degree of subvarieties
on abelian varieties}, arXiv:2606.31894v1 (2026).
\bibitem[Fu]{Fu} W. Fulton, \emph{Intersection theory}, 2nd ed., Springer,
1998.
\bibitem[GFS]{GFS} O. de Gaay Fortman, S. Schreieder, \emph{Abelian
varieties with no power isogenous to a Jacobian}, Compos. Math. 161 (2025),
1404--1457; arXiv:2401.06577.
\bibitem[Mi]{Mi} J. S. Milne, \emph{Jacobian varieties}, notes (2021 version).
\bibitem[Mu]{Mu} D. Mumford, \emph{Abelian varieties}, Oxford University
Press, 1970.
\bibitem[Oo]{Oo} F. Oort, \emph{Finite group schemes, local moduli for
abelian varieties, and lifting problems}, Compos. Math. 23 (1971), 265--296.
\bibitem[OU]{OU} F. Oort, K. Ueno, \emph{Principally polarized abelian
varieties of dimension two or three are Jacobian varieties}, J. Fac. Sci.
Univ. Tokyo Sect. IA Math. 20 (1973), 377--381.
\bibitem[SS]{SS} F. Scavia, F. Suzuki, \emph{On direct summands of products
of Jacobians over arbitrary fields}, arXiv:2507.09821v2 (2025).
\bibitem[Vo]{Vo} C. Voisin, \emph{Cycle classes on abelian varieties and the
geometry of the Abel--Jacobi map}, Pure Appl. Math. Q. 20 (2024), 2469--2496.
\end{thebibliography}

\end{document}
